% =====================================================================
%  PSAMR at fourth order in space and third order in time
%
%  Design note and staged plan, written before any code change.  The
%  predictions and falsifiers of S8 are fixed here; results go to S10 as
%  they are measured, refuted predictions recorded as such.  The prose
%  companion documents are docs/MODEL_AND_ALGORITHM.md (the scheme as it
%  stands) and docs/DESIGN_DECISIONS.md (the refuted list, which binds
%  this plan too).  File names only; line numbers go stale.
%
%  Build:  pdflatex PLAN_high_order.tex   (twice, for references)
% =====================================================================
\documentclass[11pt,a4paper]{article}

\usepackage[margin=25mm]{geometry}
\usepackage{amsmath,amssymb}
\usepackage{booktabs}
\usepackage{array}
\usepackage{enumitem}
\usepackage[colorlinks=true,linkcolor=blue,citecolor=blue,urlcolor=blue]{hyperref}
\usepackage{xcolor}
\usepackage{url}
\urlstyle{tt}
\tolerance=2000
\emergencystretch=2em

\newcommand{\dd}[2]{\frac{\partial #1}{\partial #2}}
\newcommand{\code}[1]{\texttt{\small #1}}
\DeclareUrlCommand\src{\def\UrlFont{\color{gray}\small\ttfamily}}
\newcommand{\Dt}{\Delta t}
\newcommand{\Dx}{\Delta x}
\newcommand{\Ap}{\mathcal{A}^{+}}
\newcommand{\Am}{\mathcal{A}^{-}}

\title{PSAMR at fourth order in space and third order in time\\
       \large Design note, staging and falsifiers}
\author{Branch \code{high-order} off \code{main} at \code{e9eecf18}}
\date{15 November 2026, written before any run}

\begin{document}
\maketitle

\begin{abstract}
\noindent
This note fixes the route by which the Pelanti--Shyue path of PSAMR is to
reach fourth-order face reconstruction in space and third-order,
strong-stability-preserving time integration, with the adaptive mesh
coupling consistent at the same orders.  It states what the code already
has, the options considered and the one chosen, the equations of the chosen
scheme, the two places where the achievable order is bounded by something
other than the scheme, the GPU design point, and a staged plan in which
every stage carries its predictions and falsifiers before it is run.  The
claim at the end of the plan is stated precisely in \S\ref{sec:claim},
because ``fourth order'' means different things in one and in several
dimensions.
\end{abstract}

\tableofcontents

% =====================================================================
\section{The claim this plan makes, and the one it does not}
\label{sec:claim}
% =====================================================================

The plan delivers, on the Pelanti--Shyue path (\code{PSAMR.ps\_hydro=1},
\code{ps\_flux=wp}):

\begin{enumerate}[nosep]
\item face states built by a fourth-order-accurate reconstruction of the
      cell averages (WENO5-Z, dimension by dimension, on characteristic
      fields), in place of the raw cell averages used now;
\item the Shu--Osher three-stage strong-stability-preserving Runge--Kutta
      method, SSPRK(3,3), in place of the single-step scheme (\code{do\_mol=0})
      and the Heun method (\code{do\_mol=1});
\item coarse-to-fine interpolation at fourth order in space and
      ghost cells at third order in time across coarse--fine boundaries
      under subcycling, with reflux carrying the Runge--Kutta stage
      weights.
\end{enumerate}

Measured, this means: the design rate on the smooth two-dimensional
advection test (\code{Exec/CO2\_ADV2D}) at or above four in space when the
time step is driven down faster than the mesh, and three in time on a fixed
grid; the one-dimensional battery passing its acceptance gate with errors
against the exact references no worse than now; and the two-dimensional
cases running to their stop times with the guards quiet.

It does not deliver formal fourth order on smooth problems in two and three
dimensions with nonlinear fluxes.  That needs the McCorquodale--Colella
conversions between cell averages and point values (the $h^2/24$ Laplacian
corrections on the states before the equation of state is queried, on the
face fluxes, and on the sources).  Those touch \src{PS_ctoprim.H}, the
relaxation operator and every source, which act on cell averages, and they
are parked (\S\ref{sec:parked}).  For a flashing jet with shocks and a
phase interface the difference is not measurable; the choice is recorded so
that nobody reads the ADV2D rate as a claim about the pipe-break.

Time integration is bounded separately by the operator splitting of the
relaxation sources (\S\ref{sec:bounds}): the hydro step becomes third
order, the full step with finite-rate relaxation is second order under
Strang splitting and first under Lie, as now.  Raising that needs an
implicit--explicit Runge--Kutta treatment of the rate terms, which is its
own plan.

% =====================================================================
\section{Where the code stands}
\label{sec:now}
% =====================================================================

\begin{itemize}[nosep]
\item \textbf{Three hydro paths.}  The single-fluid Godunov path
      (\src{Source/Hydro/Godunov}, PPM or PLM faces, characteristic
      tracing, second order in time, unsplit); the single-fluid
      method-of-lines path (\src{Source/Hydro/MOL}, PLM faces, Heun in
      \src{PSAMR_advance.cpp}, \code{cfl} capped at 0.3); and the
      Pelanti--Shyue wave-propagation path (\src{Source/Hydro/PelantiShyue},
      Berger--LeVeque fluctuation form, HLLC decomposition, one scalar
      van Leer limiter per wave on the Lax--Wendroff correction, face
      states the raw cell averages by design, directionally split in two
      and three dimensions by default, a contact-only transverse option).
      The wave-propagation path runs under either time scheme; under
      \code{do\_mol=1} the vanished-phase fold is applied to the stage-one
      state as well.
\item \textbf{AMR.}  AMReX \code{AmrLevel} with subcycling; coarse-to-fine
      interpolation by \code{cell\_cons\_interp} (limited linear, second
      order) in \src{PSAMR_setup.cpp}; fine ghost cells in time by linear
      interpolation between the coarse old and new states; conserved slots
      refluxed through the standard register; the volume fraction co-moved
      in symmetric capacity form in \code{PSAMR::reflux} (MODEL \S4.9).
      Ghost width \code{numGrow()} is 4 without embedded boundaries.
\item \textbf{Sources.}  The coupled relaxation operator
      (\code{ps\_relax\_mode=5}: mechanical equilibrium as a constraint,
      finite-rate thermal relaxation, statistical-rate-theory mass
      transfer, nucleation-only flash) applied once per step after the
      hydro (Lie, \code{ps\_strang=0}) or in half steps around it
      (Strang, \code{ps\_strang=1}).
\item \textbf{What AMReX already provides} (checked in the checkout at
      \code{../amrex}, \src{Src/Amr/AMReX_AmrLevel.H},
      \src{Src/AmrCore/AMReX_FillPatcher.H},
      \src{Src/AmrCore/AMReX_Interpolater.H}):
      \code{AmrLevel::RK(order,\dots)} for orders 2, 3 and 4, where order 3
      is SSPRK(3,3); \code{FillPatcher::fillRK}, which stores the coarse
      level's Runge--Kutta stage values (\code{storeRKCoarseData}) and
      fills fine ghost cells at each fine stage from a Taylor expansion of
      those stages, so the coarse--fine ghost data are consistent with the
      Runge--Kutta order under subcycling; and the interpolaters
      \code{quartic\_interp} (\code{CellConservativeQuartic}) and
      \code{cell\_quartic\_interp}.  The quartic interpolaters carry no
      limiter; they are built for smooth data.  Stage~0 confirms this by
      reading the source.
\item \textbf{Instruments.}  The one-dimensional battery and its
      fingerprints (\src{Exec/CO2_RiemannSuite/characterize.py},
      \code{verify\_canonical.py}, \code{err\_vs\_analytic.py},
      \code{gen\_convergence.py}); the smooth-advection order test
      (\src{Exec/CO2_ADV2D/inputs}, with \code{inputs-fb-rk2} for reflux
      conservation through a static fine box); the XC2D coarse--fine
      ray-diff (VERIFICATION \S7.3); T-Blowdown (\S7.1); the
      \code{[PS-GUARD]} counters and the \code{ps\_validate} realizability
      checks.
\end{itemize}

% =====================================================================
\section{Options considered}
\label{sec:options}
% =====================================================================

\begin{center}
\small
\begin{tabular}{p{0.30\textwidth}p{0.30\textwidth}p{0.32\textwidth}}
\toprule
Route & For & Against \\
\midrule
Method of lines: fourth-order reconstruction at faces, HLLC fluctuations,
SSPRK(3,3) \textbf{(chosen)} &
Keeps the fluctuation form, the recovered fluxes, the per-cell deposits and
the HLLC star state; AMReX supplies the time scheme and the coarse--fine
coupling; the recipe of Johnsen--Colonius and Coralic--Colonius for
non-conservative two-phase models, and of Ketcheson, Parsani and LeVeque
for high-order wave propagation. &
Three stage evaluations per step and the unsplit CFL bound; a new in-cell
term for the non-conservative slots; formally second order in several
dimensions without the point-value conversions (\S\ref{sec:claim}). \\
\addlinespace
Method of lines with the PPM interpolant of \src{Godunov/PPM.cpp} &
Lowest effort; the interpolant is in the tree; sharper on contacts than
WENO5. &
The limiter degrades to second order at every extremum; fourth order only
in the unlimited smooth regime.  Kept as the fallback reconstruction if
WENO5 fails a Stage~3 falsifier. \\
\addlinespace
ADER with path-conservative WENO (Dumbser--Toro) &
Single step, so no stage ghost fills; stiff sources inside the local
space-time predictor at full order. &
A different code: a local space-time predictor with many equation-of-state
evaluations per cell per step, on a code whose measured cost item is the
equation of state (FUTURE\_WORK \S11); not an AMReX pattern. \\
\addlinespace
Two-stage fourth-order Lax--Wendroff on a generalised Riemann solver
(Li--Du) &
Two stages for fourth order in time. &
Needs a generalised Riemann solver for the six-equation model with a
real-fluid equation of state.  Research. \\
\addlinespace
Discontinuous Galerkin; finite-difference WENO &
Formal order on any mesh (DG); cheap point-value fluxes (FD). &
DG needs subcell limiting for flashing jets and has no AMReX support; FD
WENO is not conservative at coarse--fine faces without special treatment,
and the reflux machinery is built for finite volumes. \\
\bottomrule
\end{tabular}
\end{center}

\paragraph{Why strong stability preserving.}
Each stage of an SSP Runge--Kutta method is a convex combination of forward
Euler steps.  Any property the forward Euler step preserves under its CFL
bound is therefore preserved by the full step under that bound times the
SSP coefficient.  Here the forward Euler step is the first-order
wave-propagation update that the floors, folds and guards already protect,
and the properties that matter are positivity of the partial masses,
$\alpha_1\in[0,1]$, phase energies inside the equation-of-state domain,
and the bound of the limited correction.  A third-order scheme that is not
SSP gives no such guarantee on its intermediate states, and in this code an
intermediate state that leaves the domain is a refused face or an aborted
run, counted by the guards.  SSPRK(3,3) has SSP coefficient 1, so its CFL
bound is the forward Euler bound exactly, and it runs in two registers in
Shu--Osher form.  SSPRK(4,3) does a third less work per unit of simulated
time (four stages at coefficient 2) and is the natural second step; it is
not in AMReX, so it needs its own Taylor coefficients in
\code{FillPatcher}, and it is deferred until SSPRK(3,3) is measured.  No
four-stage fourth-order SSP scheme with positive coefficients exists; the
target here is third order in time, so that barrier is not met.

% =====================================================================
\section{The scheme}
\label{sec:scheme}
% =====================================================================

\subsection{Semi-discretisation in fluctuation form}

With $Q_i$ the vector of cell averages of the Pelanti--Shyue state (MODEL
\S1) and $q^{L}_{i+1/2}, q^{R}_{i+1/2}$ the reconstructed states on the two
sides of face $i+1/2$, the method-of-lines update along one direction is
\begin{equation}
\label{eq:mol}
\frac{dQ_i}{dt} = -\frac{1}{\Dx}\Big(\Ap\Delta q_{i-1/2} + \Am\Delta q_{i+1/2}\Big)
                  -\frac{1}{\Dx}\int_{x_{i-1/2}}^{x_{i+1/2}} A\big(q_i(x)\big)\,\dd{q_i}{x}\,dx ,
\end{equation}
where $\Ap\Delta q$ and $\Am\Delta q$ are the right- and left-going
fluctuations of the HLLC decomposition of the jump $q^{R}-q^{L}$ at the
face, as now (\src{PS_hllc.H}, \src{PS_umeth.cpp}), and $q_i(x)$ is the
reconstruction polynomial inside cell $i$.  The second term is new: with
raw cell averages at the faces it vanishes, because $q_i(x)$ is constant.
On the conserved slots it is $f(q^{L}_{i+1/2})-f(q^{R}_{i-1/2})$ exactly,
so together with the fluctuations the update telescopes to the recovered
flux form that \src{Hydro_consup.cpp} already differences.  On the
non-conserved slots ($\alpha_1$, the two phase energies) it is the
integral of the non-conservative product across the cell, evaluated by
quadrature on the reconstruction polynomial with the velocity taken from
the same polynomial, and deposited per cell as the correction term is
deposited now (MODEL \S3.3).  This is the Ketcheson--Parsani--LeVeque form
of high-order wave propagation; for the $\alpha$ equation it is the
Johnsen--Colonius treatment of $u\,\partial_x\alpha$, and the face
contribution keeps the advective split
$\bar u\,\Delta\alpha$ of MODEL \S3.6, not $\Delta(\alpha u)$.

The limited Lax--Wendroff correction of the present scheme
(\code{ps\_wp\_order=2}) is not applied under the new path: its role, to
raise the order above one, is taken by the reconstruction.  The
contact-wave taper (MODEL \S3.4) has no counterpart yet; whether a
high-order contact needs one is a Stage~3 measurement.

\subsection{Reconstruction}

WENO5-Z reconstruction of the cell averages to the two face values, on
the characteristic fields of the frozen six-equation system projected at
each face with the Roe-averaged or arithmetic-mean state (two acoustic
fields $u\pm c$, five contact fields at $u$; the eigenvectors of
\src{PS_wavespeed.H} and \src{PS_hllc.H}).  Characteristic-wise
reconstruction is the high-order analogue of the one-scalar-per-wave
limiter: the DESIGN\_DECISIONS record that any per-slot limiting bends the
wave (the rejected $\phi=\min(\phi_1,\phi_2)$ entry, MODEL \S3.3) binds
here, and reconstructing the slots independently is not an option.  The
variables reconstructed are $\alpha_1$, $m_1$, $m_2$, the momenta and the
two phase energies; the mixture slots are formed as sums at the face so the
identities $m_1+m_2=\rho$ and $\mathcal{E}_1+\mathcal{E}_2=\rho E$ hold at
the face by construction.  Five cells enter each face value; the ghost
width of 4 covers it.  Reconstruction near a vanished or corridor phase is
the delicate point and the reason the present scheme uses raw averages:
the presence classifier (\src{PS_presence.H}) runs on the reconstructed
face states, and a face whose reconstruction leaves a present phase's
$\alpha_k$ or $m_k$ outside the admissible set falls back to the raw
averages at that face and is counted (\code{recon\_fallback} in
\src{PS_counters.H}).  Whether that fallback fires anywhere in the
battery or the two-dimensional cases is a Stage~3 prediction.

The PPM interpolant of \src{Godunov/PPM.cpp} is the fallback reconstruction
if WENO5-Z fails a Stage~3 falsifier, with the same characteristic
projection.

\subsection{Time integration}

SSPRK(3,3) in Shu--Osher form, with $L$ the right-hand side of
\eqref{eq:mol} summed over directions and $\Pi$ the per-stage
post-processing described below:
\begin{align}
Q^{(1)} &= \Pi\big[Q^n + \Dt\,L(Q^n)\big], \\
Q^{(2)} &= \Pi\big[\tfrac34 Q^n + \tfrac14\big(Q^{(1)} + \Dt\,L(Q^{(1)})\big)\big], \\
Q^{n+1} &= \Pi\big[\tfrac13 Q^n + \tfrac23\big(Q^{(2)} + \Dt\,L(Q^{(2)})\big)\big].
\end{align}
$\Pi$ is the vanished-phase fold already applied to the stage-one state
under \code{do\_mol=1}, the mixture-mass resync, and the mechanical
equilibrium projection of the relaxation operator (the constraint part of
mode 5, applied without its finite-rate parts).  Each is a map of
admissible states to admissible states, so each stage remains a convex
combination of admissible forward Euler steps and the SSP property is
kept.  The finite-rate thermal relaxation, mass transfer and the flash
stay outside the stages: Strang halves around the three-stage step under
\code{ps\_strang=1}, one Lie step after it under \code{ps\_strang=0}.

The time step keeps the per-direction definition of \src{Timestep.H},
$\Dt = \code{cfl}\cdot\min_d \Dx_d/\lambda_d$.  The forward Euler bound of
the unsplit operator \eqref{eq:mol} in that definition is
$\code{cfl}\le 1/\mathrm{DIM}$; the present MOL cap of 0.3 in
\src{PSAMR.cpp} is kept as the default and its relaxation to $1/\mathrm{DIM}$
is a Stage~1 measurement, not a change made on trust.  The battery runs
at \code{cfl=0.25}.  Under the new scheme \code{ps\_wp\_transverse} must be
0: the method of lines is unsplit by construction and the transverse term
belongs to the split single-step form.

\subsection{AMR coupling}

\begin{itemize}[nosep]
\item \textbf{Time.}  The level advance calls \code{AmrLevel::RK(3,\dots)}
      instead of the hand-written Heun.  The coarse level stores its three
      stage values; each fine stage fills its ghost cells through
      \code{FillRKPatch}, which evaluates the coarse Taylor expansion at
      the fine stage time.  The hand-written Heun path remains as
      \code{PSAMR.rk\_order=2}, the default until the record is in.
\item \textbf{Space.}  \code{quartic\_interp} replaces
      \code{cell\_cons\_interp} for the state under
      \code{PSAMR.cf\_interp=quartic} (default \code{cons}, the present
      behaviour).  It is unlimited; the post-interpolation presence guards
      of \src{PSAMR.cpp} (the ``post-C-F-interp'' range print and the
      sliver handling after regrid) are the instrument that says whether a
      limited fourth-order interpolation is needed, and Stage~2 decides on
      that measurement.
\item \textbf{Reflux.}  The flux register accumulates the stage fluxes with
      the Shu--Osher weights ($\tfrac16,\tfrac16,\tfrac23$ on the
      stage fluxes in Butcher form), so the refluxed flux is the one the
      step used.  The volume-fraction co-move is unchanged in form: it acts
      on the refluxed masses at the pre-reflux phase densities, whatever
      produced them.
\item \textbf{Average-down} is exact on cell averages and does not change.
\item \textbf{Boundaries.}  The NSCBC outflow ghosts (\src{PS_nscbc.H}) are
      built per stage from the stage state.  Embedded-boundary builds are
      out of scope: the CO$_2$ decks build with \code{USE\_EB=FALSE}, and a
      cut cell would stay second order whatever the interior does.
\end{itemize}

% =====================================================================
\section{What bounds the order}
\label{sec:bounds}
% =====================================================================

\begin{enumerate}[nosep]
\item \textbf{The relaxation split.}  With the rate terms outside the
      stages the full step is second order in time at best.  An
      implicit--explicit scheme (Pareschi--Russo, or ARS(4,4,3)) with the
      rate terms implicit would give third order and asymptotic
      preservation; the implicit SSP coefficients of such schemes are small,
      so what they buy is positivity and the correct stiff limit, not an
      SSP bound.  The mechanical constraint is a projection and goes inside
      the stages at no cost to the order.  The flash is event-like and is
      low order by nature.  Stage~4 measures the temporal order of the full
      step and records it; the implicit--explicit treatment is its own plan.
\item \textbf{Averages against point values.}  \S\ref{sec:claim}.  For
      the constant-velocity advection of ADV2D the dimension-by-dimension
      scheme is exact at the design order, because averaging commutes with
      a linear flux; that is why ADV2D can measure the reconstruction order
      while the pipe-break cannot.
\item \textbf{Shocks and the interface.}  Behind a shock the observed order
      is one, as it is now (VERIFICATION \S6.4), and the limiter in the
      WENO weights reduces the rate at extrema.  The battery measures
      errors against exact references, not rates, and the prediction is
      ``no worse'', not ``fourth order''.
\end{enumerate}

% =====================================================================
\section{The GPU design point}
\label{sec:gpu}
% =====================================================================

No new mechanism.  Each stage is the present set of \code{ParallelFor}
launches (reconstruction, face decomposition, deposit, consup) once more,
so launches per step triple and the equation-of-state calls per cell per
step triple with them; the two-register Shu--Osher form adds one state-sized
\code{MultiFab} over the Heun path; the coarse level holds three stage
\code{MultiFab}s per state for \code{fillRK}.  The characteristic
projection is a per-face $7\times7$ product in registers.  The new counter
is host-only and compiles to a no-op on the device, as \code{flux\_sanit}
does.  Predicted cost per unit of simulated time against the present split
single-step scheme at \code{cfl=0.9}: three stages at \code{cfl=0.3}
against one stage at 0.9 gives a factor of about 9 on the hydro, less on
the whole step because the relaxation operator runs once per step either
way; with the cap relaxed to $1/\mathrm{DIM}$ in two dimensions, about 5.
The equation-of-state call count is the thing to watch (FUTURE\_WORK
\S11).  GPU runs are not bitwise reproducible run to run, so GPU
predictions are stated to round-off, never bitwise.

% =====================================================================
\section{Dials}
\label{sec:dials}
% =====================================================================

All new keys default to the present behaviour, so an unchanged deck runs
unchanged.
\begin{center}
\small
\begin{tabular}{llp{0.55\textwidth}}
\toprule
Key & Default & Meaning \\
\midrule
\code{PSAMR.rk\_order} & 2 & Under \code{do\_mol=1}: 2 = the hand-written Heun step; 3 = \code{AmrLevel::RK(3)}, SSPRK(3,3) with \code{fillRK} ghosts. Any other value aborts. \\
\code{PSAMR.cf\_interp} & \code{cons} & Coarse-to-fine interpolater for the state: \code{cons} = \code{cell\_cons\_interp}; \code{quartic} = \code{quartic\_interp}. \\
\code{PSAMR.ps\_wp\_order} & 2 & Extended: 4 = reconstructed faces (WENO5-Z characteristic-wise) with the in-cell term and no Lax--Wendroff correction; 1 and 2 as now. \\
\code{PSAMR.ps\_wp\_recon} & \code{weno5z} & Under \code{ps\_wp\_order=4}: \code{weno5z} or \code{ppm}. \\
\bottomrule
\end{tabular}
\end{center}
\code{ps\_wp\_order=4} with \code{ps\_wp\_transverse} nonzero aborts.
\code{ps\_wp\_order=4} with \code{do\_mol=0} aborts: the reconstructed
faces are a method-of-lines operator and have no characteristic-tracing
half step.

% =====================================================================
\section{Staging, with predictions and falsifiers}
\label{sec:stages}
% =====================================================================

Every stage builds on the Mac (\code{COMP=llvm}) and, before it is called
done, on peart (\code{COMP=gnu}, CUDA).  Comparisons use
\path{~/src/PSAMR-trial/bench_nov14/pfdiff.py} (per level,
max $|$difference$|$/range per field).  Control binaries are copied aside
before each rebuild.  Every two-dimensional run, new and control, carries
\code{PSAMR.ps\_diag\_mass=1} so the \code{[PS-GUARD]} line prints.  The
battery is run with \code{EXE=} named explicitly.  Where a prediction
fails, the failure is recorded in \S\ref{sec:results} before anything is
changed.

\subsection*{Stage 0: controls, no code change}

\begin{enumerate}[nosep]
\item Read \code{CellConservativeQuartic} in the AMReX checkout and record
      whether it limits.  Prediction: it does not.
\item ADV2D at $32^2, 64^2, 128^2$, \code{do\_mol=0}, \code{ps\_wp\_limiter=none}:
      record $L_1(\alpha)$ and the rate.  Prediction: rate $2.0\pm0.1$,
      $L_1=6.80\times10^{-4}$ at $64^2$ as VERIFICATION \S7.2 records.
\item ADV2D at $64^2$ under \code{do\_mol=1} (Heun, \code{ps\_wp\_order=2})
      at \code{cfl} 0.3, 0.15, 0.075: record the self-convergence ratio in
      $\Dt$.  Prediction: $\log_2$ of the ratio $2.0\pm0.2$ (Heun is second
      order).  Falsifier: below 1.8, which would mean the Heun path is not
      the second-order step it is documented as, and Stage~1 compares
      against a wrong control.
\item Battery under \code{do\_mol=1} Heun: record the \code{verify\_canonical}
      PASS set and the \code{err\_vs\_analytic} errors.  These are the
      Stage~1 and Stage~3 controls; no prediction, a record.
\item Wall time per step and \code{[PS-GUARD]} counters for XC2D (40
      steps) and T-Blowdown (\code{inputs-x-amr}, 300 steps) under
      \code{do\_mol=0} and under \code{do\_mol=1} Heun.  Record.
\end{enumerate}

\subsection*{Stage 1: SSPRK(3,3) through \code{AmrLevel::RK}, spatial operator unchanged}

\code{rk\_order=3}, \code{ps\_wp\_order=2}, \code{cf\_interp=cons}.  The
hydro right-hand side is the present one; only the time scheme and the
ghost filling change.

\begin{enumerate}[nosep]
\item ADV2D $64^2$, \code{cfl} 0.3, 0.15, 0.075: $\log_2$ of the
      self-convergence ratio in $\Dt$ is $3.0\pm0.3$.  Falsifier: below 2.5
      (then the stage weights or the ghost filling are wrong, and no
      further stage runs until it is found).
\item ADV2D $32^2$ to $128^2$ at fixed \code{cfl}: rate $2.0\pm0.1$, and
      $L_1$ at $64^2$ within 5\,\% of the Heun value (the spatial error
      dominates at second order).  Falsifier: rate below 1.9.
\item ADV2D \code{inputs-fb-rk2} (static fine box, periodic): total mass and
      $\rho E$ conserved to round-off across the \code{sum\_interval=20}
      totals (below $10^{-13}$ relative).  Falsifier: a drift above that,
      which means the reflux register is not carrying the stage weights
      the step used.
\item Battery: the \code{verify\_canonical} PASS set equal to the Heun
      control's; every \code{err\_vs\_analytic} error within 2\,\% of the
      Heun control's (the battery is shock-limited and the spatial operator
      is unchanged).  Falsifier: a check that passed under Heun and fails
      here, or an error moved by more than 10\,\%.
\item XC2D, 40 steps, level 1: runs to step 40; \code{flux\_sanit} and the
      three limiter counters 0; the \code{reflux\_*} counters equal to the
      Heun control's.  Falsifier: an abort, or a limiter counter above 0.
\item T-Blowdown 300 steps: runs; counters 0; final plotfile against the
      Heun control within $10^{-3}$ of range on every field (a different
      second-order-in-time error, not a different solution).  Falsifier:
      above $10^{-2}$ anywhere.
\item CFL cap: ADV2D and XC2D at \code{cfl=0.45} under \code{rk\_order=3}:
      stable, counters 0.  Prediction: stable (the forward Euler bound in
      two dimensions is 0.5).  Falsifier: an abort or a counter above 0,
      in which case the cap stays at 0.3 and the cost prediction of
      \S\ref{sec:gpu} is the factor 9, not 5.
\item GPU (peart): XC2D 40 steps against the Mac run, every field to
      round-off (below $10^{-12}$ of range on both levels).
\end{enumerate}

\subsection*{Stage 2: fourth-order coarse-to-fine interpolation}

\code{cf\_interp=quartic} on top of Stage 1.

\begin{enumerate}[nosep]
\item ADV2D \code{inputs-fb-rk2}: conservation to round-off as in Stage 1
      (the quartic interpolater is conservative).  The $L_1(\alpha)$
      against the translated initial condition at the stop time is no
      larger than under \code{cons}.  Falsifier: larger by more than 10\,\%.
\item XC2D 40 steps, the coarse--fine ray-diff of VERIFICATION \S7.3
      against the uniform reference: the contact gap $|\Delta\alpha_1|$
      (now $8.5\times10^{-8}$) and the relative phase-energy gap in the
      coarse--fine band (now $1.9\times10^{-4}$) do not grow.  Prediction:
      both fall.  Falsifier: either grows by more than a factor 2.
\item The post-C-F-interp range print of \src{PSAMR.cpp} on XC2D and
      T-Blowdown: $\alpha_1$ after interpolation stays in $[0,1]$ and the
      sliver counters after regrid equal the \code{cons} control's.
      Prediction: the unlimited quartic interpolant produces no
      out-of-range $\alpha_1$ on these two cases, because both carry a
      resolved contact.  Falsifier: an out-of-range value, which decides
      that a limited fourth-order interpolant (the McCorquodale--Colella
      limited quartic) is needed before Stage 3, and that is recorded as a
      refuted prediction and a new item.
\item Pipe-break \code{inputs.satjet\_demo2} to step 200 (the flashing jet
      is where an unlimited interpolant will show if it shows anywhere):
      runs; the post-C-F-interp print never reports $\alpha_1$ outside
      $[0,1]$; \code{reflux\_*} counters within a factor 2 of the control's.
      Falsifier: an out-of-range $\alpha_1$, same decision as item 3.
\item GPU: XC2D to round-off against the Mac, as in Stage 1.
\end{enumerate}

\subsection*{Stage 3: reconstructed faces and the in-cell term}

\code{ps\_wp\_order=4}, \code{ps\_wp\_recon=weno5z}, on top of Stage 2.

\begin{enumerate}[nosep]
\item ADV2D space: $32^2, 64^2, 128^2$ with \code{cfl} scaled as
      $(32/N)^{2/3}$ (so $\Dt\propto h^{5/3}$ and the third-order temporal
      error falls faster than $h^4$), \code{ps\_wp\_limiter=none} has no
      meaning here and is not set: rate at or above $4.0$ in
      $L_1(\alpha)$.  Prediction: $4.5\pm0.5$ (WENO5-Z is fifth order on
      smooth data; the nonlinear weights cost a little).  Falsifier: below
      3.5, in which case the reconstruction, the characteristic projection
      or the in-cell term is wrong, and the PPM fallback
      (\code{ps\_wp\_recon=ppm}) is run to tell which: PPM at or above 3.5
      with WENO below it points at the WENO weights; both below it point
      at the in-cell term or the projection.
\item ADV2D time: $64^2$, \code{cfl} 0.3, 0.15, 0.075: $\log_2$ ratio
      $3.0\pm0.3$ as in Stage 1 (the spatial error is now small enough
      that the temporal one shows).  Falsifier: below 2.5.
\item ADV2D conservation (\code{inputs-fb-rk2}) to round-off, as before.
\item The in-cell term in one dimension: a uniform-velocity, uniform-pressure
      advected contact in $\alpha_1$ (battery case B11's configuration with
      relaxation off, \code{ps\_do\_relax=0}) keeps $u$ and $P$ uniform to
      round-off ($10^{-12}$ relative) for the whole run.  This is the
      Abgrall--Saurel condition; a wrong in-cell term for $\alpha$ or the
      phase energies breaks it in the first step.  Falsifier: a pressure or
      velocity disturbance above $10^{-10}$ at the contact.
\item Battery: \code{verify\_canonical} PASS set equal to the Stage~1
      set; every \code{err\_vs\_analytic} error no larger than the Stage~1
      value, and the errors on the smooth parts of the fans smaller.
      Realizability (\code{ps\_validate} V2, V3, V6, V7) with zero
      violations; \code{energyid} on B9 no larger than now
      ($2.9\times10^{-3}$); wall reflection on B8 and B12 mirror-exact to
      round-off (the characteristic projection and WENO weights are
      symmetric under $x\to-x$ if written so; a broken mirror is a coding
      error, not a scheme property).  \code{recon\_fallback} 0 on every
      case except those with a vanished phase present (B7, B9), where it is
      recorded, not predicted.  Falsifiers: a check lost, an error larger
      than Stage~1, a realizability violation, a broken mirror.
\item Self-convergence (\code{gen\_convergence.py}, N = 128, 256, 512): the
      observed order on the shocked cases stays near 1 and on the smooth
      cases rises above 2.  Record; no hard falsifier, the battery errors
      are the gate.
\item XC2D 40 steps: runs; the coarse--fine gaps of Stage 2 do not grow;
      counters as Stage 2 plus \code{recon\_fallback}, recorded.
      T-Blowdown 300 steps: runs; counters 0 or recorded; the orientation
      suite of VERIFICATION \S7.1 still agrees across orientations to
      round-off.  Falsifier: an abort, or an orientation asymmetry.
\item Pipe-break to step 200: runs; centreline symmetry to round-off; the
      near-orifice shear-mode checkerboard of MODEL \S3.8 is recorded
      (amplitude at the same station as the record), with no prediction
      that it is smaller.  Falsifier: an abort, or broken centreline
      symmetry.
\item GPU: battery fingerprints are host-only; XC2D and T-Blowdown to
      round-off against the Mac.  Wall time per step on peart against the
      Stage~0 control: within a factor 4 of \code{do\_mol=1} Heun per step
      (three stages against two, plus the reconstruction).  Falsifier:
      above 6, which means a launch-count or equation-of-state call
      problem to find before the stage is done.
\end{enumerate}

\subsection*{Stage 4: the relaxation placement}

On top of Stage 3.  Mechanical projection inside the stages; rate terms
outside under \code{ps\_strang=1}.

\begin{enumerate}[nosep]
\item A smooth finite-rate case (B11 with \code{ps\_relax\_mode=2}, the
      finite-rate control closure, at a $\tau$ comparable to the step): the
      $\log_2$ self-convergence ratio in $\Dt$ on a fixed grid is
      $2.0\pm0.3$ under Strang and $1.0\pm0.3$ under Lie.  This is the
      record of the bound of \S\ref{sec:bounds}, item 1.  Falsifier for
      the scheme: Strang below 1.7, which would mean the projection inside
      the stages interferes with the halves (then the projection is moved
      out and the measurement repeated).
\item Battery under mode 5 with Strang against mode 5 with Lie: the PASS
      set equal; errors recorded.
\item Pipe-break to step 200 under Strang: runs; the mode-5 evidence
      quantities of VERIFICATION \S7.6 recorded at the same stations.
\item GPU as before.
\end{enumerate}

\subsection*{Stage 5: the record}

\code{VERIFICATION.md} \S7.2 gains the spatial and temporal rates and their
run lines; \S6 gains the Abgrall--Saurel contact check; \code{MODEL\_AND\_ALGORITHM.md}
\S3 gains the reconstructed-face operator, the in-cell term and the
characteristic projection, with the raw-average scheme kept as
\code{ps\_wp\_order=2}; \code{DESIGN\_DECISIONS.md} gains one entry per
decision taken here (SSP over non-SSP; characteristic-wise over slot-wise
reconstruction; the unlimited quartic interpolant kept or replaced on the
Stage~2 measurement; the CFL cap); \code{RUNNING.md} gains the four dials.
Then the defaults are decided by Marc on the record, not here.

% =====================================================================
\section{Parked}
\label{sec:parked}
% =====================================================================

\begin{itemize}[nosep]
\item The McCorquodale--Colella average-to-point conversions for formal
      multi-dimensional fourth order, including the equation of state and
      the sources at point values.
\item The implicit--explicit Runge--Kutta treatment of the finite-rate
      relaxation terms, for third order in the full step.
\item SSPRK(4,3), with its own \code{FillPatcher} coefficients, for a third
      less work per simulated time.
\item A limited fourth-order coarse-to-fine interpolant, if Stage~2 shows
      the unlimited one failing.
\item Interface sharpening (THINC-type) for the contact, if the WENO5
      contact is measurably more smeared than the present tapered scheme on
      the pipe-break.
\item Embedded boundaries at any order above two.
\end{itemize}

% =====================================================================
\section{Commits}
\label{sec:commits}
% =====================================================================

One item per commit, on \code{high-order} off \code{main} at
\code{e9eecf18}; nothing pushed without Marc's word.

\begin{enumerate}[nosep]
\item This plan, and its row in \code{docs/README.md}.
\item Stage 0 controls appended to \S\ref{sec:results}.
\item Stage 1 code: \code{PSAMR.rk\_order}, the \code{AmrLevel::RK(3)}
      call in \src{PSAMR_advance.cpp}, the stage-weighted flux register,
      \code{RUNNING.md} row; then its record.
\item Stage 2 code: \code{PSAMR.cf\_interp} in \src{PSAMR_setup.cpp};
      then its record.
\item Stage 3 code: the reconstruction, the projection and the in-cell term
      in \src{PS_umeth.cpp} and a new \src{PS_recon.H}, the counter,
      the dials; then its record.
\item Stage 4: the projection placement in the stages; then its record.
\item Stage 5: the documents.
\end{enumerate}

% =====================================================================
\section{Results}
\label{sec:results}
% =====================================================================

Empty at the time of writing.  Each stage appends its measured numbers
against the predictions above, with refuted predictions marked as such.

\end{document}
